% =============================================================================
% INTRODUCTION
% Chapter-level intro — NOT a numbered chapter
% =============================================================================
% COMPLETELY REWRITTEN — Old Knowledge Graph / Data Fabric content removed
% Project: A Context-Aware Framework for Streaming Data Quality Monitoring
% Date: April 24, 2026
% =============================================================================

\section{Motivation and Problem Statement}
\label{sec:intro-motivation}

Modern transportation systems — including urban bus fleets, taxi services, and
rail networks — generate continuous streams of Global Positioning System (GPS)
position records at update intervals ranging from seconds to minutes. These records
constitute the primary data source for real-time operations such as arrival-time
prediction, fleet management, and service scheduling. The quality of these records
directly determines the reliability of downstream applications; erroneous or
anomalous GPS data can propagate through the entire operational pipeline,
degrading predictive accuracy and compromising service quality.

GPS position records in public transit systems are susceptible to multiple classes of
data quality defects. Sensor noise produces implausible speed measurements;
GPS spoofing can inject fabricated vehicle positions separated by hundreds of meters
within seconds; transmission errors cause duplicate event records; and stale
position updates may propagate outdated location information to downstream
applications. These defects are not merely theoretical concerns — the NYC MTA Bus
system, for example, operates hundreds of vehicles updating at approximately
30-second intervals across the five boroughs of New York City, generating
millions of records per day that must be continuously monitored for quality.

Existing data quality frameworks were designed primarily for batch processing
environments. Great Expectations \citep{great_expectations} and Soda Core
\citep{soda-core} provide sophisticated constraint specification and validation
pipelines, but both operate in batch mode: data must be written to a warehouse
before quality checks can execute. In batch processing, anomalies are detected
hours after they occur, making these tools unsuitable for real-time transit
monitoring where immediate detection is operationally critical. Great Expectations
and Soda Core additionally lack cross-record validation capabilities that would
enable detection of GPS trajectory anomalies such as impossible speed jumps.

Stream DaQ \citep{streamdaq2025} represents the most closely related prior work,
providing a streaming-native data quality framework with real-time violation
detection. However, Stream DaQ was designed for generic data streams and lacks
GPS trajectory validation capabilities; it cannot detect physically implausible
vehicle speeds or coordinate jumps in transit GPS data. METER \citep{meter2024}
addresses concept drift in streaming machine learning models but does not implement
data quality validation for trajectory data. GTFS Validator \citep{gtfs_validator}
validates static GTFS schedule data but operates in batch mode and does not
support real-time GTFS-realtime feeds.

The gap we identify is the absence of a streaming data quality framework that
implements cross-record GPS trajectory validation on real-time GTFS-realtime
feeds. No surveyed framework — among Great Expectations, Soda Core, Stream DaQ,
METER, Ada-Context \citep{adacontext2025}, Deequ \citep{deequ2018}, GTFS Validator,
or any of 14 others surveyed in our literature review — implements Haversine-based
GPS speed and jump detection on streaming vehicle positions with context-aware
threshold adaptation. This gap is operationally significant: erroneous GPS data
in public transit directly degrades arrival-time predictions and fleet management
decisions that affect thousands of passengers daily.

We present \textbf{StreamDQ}: a context-aware framework for streaming data quality
monitoring that validates GPS trajectory quality on real GTFS-realtime feeds in
real time. StreamDQ is not a production-grade system; it is a research and
education platform designed to demonstrate that streaming GPS validation is feasible,
measurable, and more effective than batch tools for real-time transit data quality.

\section{Research Objectives}
\label{sec:intro-objectives}

We formulate six research questions that address the core claims of this work.
Each question is stated with a measurable hypothesis and a target threshold;
all thresholds are labeled \texttt{[TIER-2\ ESTIMATED]} because they require
benchmark validation.

\textbf{RQ1.} Do hierarchical context-aware thresholds (L0--L4) improve detection F1
over static global thresholds? We hypothesize a target improvement of
$\Delta\text{F1} \geq 5$ percentage points for events falling into L0 context
cells (fine-grained hour--zone--weekday), assessed via Wilcoxon signed-rank test
with Holm--Bonferroni correction ($\alpha_{\text{adj}} = 0.0083$).
\texttt{[TIER-2\ ESTIMATED]}

\textbf{RQ2.} Does the L4 global fallback reduce threshold specificity
detectably compared to L0 context-specific thresholds? We assess whether the
false-negative rate at L4 is higher than at L0 by more than 10\%,
measured by Friedman's test across context cells.
\texttt{[TIER-2\ ESTIMATED]}

\textbf{RQ3.} Does Trajectory Quality Scoring (TQS) correlate with a
downstream quality signal? We measure Pearson correlation $\rho > 0.70$ between
TQS composite scores and independent downstream prediction error. Critically,
we redesign RQ3 to avoid circularity: TQS is computed from raw event properties
alone (not from rule violation counts), and correlation is measured against
ETA prediction error on NYC MTA Bus, not against injection rate.
\texttt{[TIER-2\ ESTIMATED]}

\textbf{RQ4.} Does the cross-record (CRS) rule layer provide incremental F1
beyond what syntactic and semantic rules achieve alone? We target
$\Delta\text{F1} \geq 5$ percentage points from adding CRS rules to the SYN+SEM
baseline, measured via Wilcoxon signed-rank test.
\texttt{[TIER-2\ ESTIMATED]}

\textbf{RQ5.} Do CRS rules achieve precision $> 0.70$ on synthetically injected
GPS anomalies in the NYC MTA Bus GTFS-realtime stream? We assess both CRS001
(GPS speed bounds) and CRS002 (GPS jump detection) using Wilson score confidence
intervals at 95\% confidence level.
\texttt{[TIER-2\ ESTIMATED]}

\textbf{RQ6.} Does ML-augmented threshold calibration (Bayesian Optimization)
improve F1 over rule-only thresholds? This is a measured contribution:
positive, negative, and mixed results are all scientifically valuable and will
be reported as such. Phase 3 implementation is mandatory.
\texttt{[TIER-2\ ESTIMATED]}

\section{Contributions}
\label{sec:intro-contributions}

We implement four primary contributions for the StreamDQ research platform.

\textbf{C1: First streaming GPS trajectory validation on real GTFS-realtime
feeds.} We implement CRS001 (Haversine-based GPS speed bounds $[2.0,\ 100.0]$
km/h) and CRS002 (GPS jump detection $>400$ m in 30 s) as Java
\texttt{KeyedProcessFunction} operators in Apache Flink, evaluated on the live
NYC MTA Bus GTFS-realtime public feed. No prior framework implements cross-record
GPS validation on streaming transit data.
\texttt{[TIER-2\ ESTIMATED]}

\textbf{C2: Hierarchical context-aware threshold fallback (L0--L5).} We design
a five-level hierarchical threshold system that traverses from the finest-grained
context cell (L0: hour $\times$ zone $\times$ weekday/weekend, $\geq 100$
samples) to physics priors (L5), gracefully degrading through intermediate levels
(L1--L4). We validate that this fallback mechanism is absent from all
surveyed frameworks.
\texttt{[TIER-2\ ESTIMATED]}

\textbf{C3: Ground-truth evaluation methodology with explicit limitations.}
We implement a complete evaluation pipeline comprising synthetic anomaly injection,
per-entity ground-truth tracking in PostgreSQL, and precision/recall/F1 metrics
with 95\% bootstrap confidence intervals (1,000 iterations). Uniquely, we
explicitly label what cannot be measured: CRS003 recall is blocked by a known
evaluation bug (NG-4), and distributed Flink latency is not benchmarked.
\texttt{[TIER-1\ VERIFIED\ for\ methodology;\ TIER-2\ ESTIMATED\ for\ results]}

\textbf{C4: ML-augmented threshold calibration.} We design a three-component
ML pipeline — Bayesian Optimization (Priority 1, GO), Isolation Forest (Priority 2,
conditional), and XGBoost (Priority 3, conditional) — that calibrates context
thresholds without overriding rule-based violation detection. LSTM is explicitly
excluded (Priority 4, NO-GO) due to unconfirmed training data and high
redundancy with CRS002.
\texttt{[TIER-2\ ESTIMATED]}

\section{Thesis Structure}
\label{sec:intro-structure}

The remainder of this thesis is organized as follows.

\textbf{Chapter 1} presents the literature review and theoretical background,
covering streaming data quality definitions and frameworks, context-aware data
quality methods, GPS trajectory validation for transportation systems, and the
role of machine learning in data quality calibration.

\textbf{Chapter 2} describes the StreamDQ system architecture and algorithm
design, including the three-layer rule taxonomy (syntactic, semantic, and
cross-record), the hierarchical context-aware threshold system, the Trajectory
Quality Scoring mechanism, and the ML augmentation pipeline.

\textbf{Chapter 3} details the experimental evaluation methodology, including
datasets, synthetic anomaly injection, ground-truth tracking, statistical tests,
and the ablation study design. Result sections contain placeholder tables
labeled \texttt{[TIER-2\ ESTIMATED]} pending benchmark execution. A mandatory
limitations section explicitly acknowledges seven documented constraints.

The \textbf{Conclusion} summarizes the contributions and key findings, and
identifies five directions for future work, including fixing the CRS003
evaluation blocker, implementing the external context stub, and deploying a
distributed Flink cluster for full-scale benchmarks.
